
\section{Test 4: Prova de síntesi}

\iniciTest{tD}{a,c,b,d,b,b,b,c,b,b,c,a}

\begin{enumerate}
\pregunta Siguin
\[
	\left[\frac{a}{b}\right]
	=
	\left[\frac{a'}{b'}\right]
\]
i sigui \(\left[\frac{c}{d}\right]\in\Q\). Quina igualtat expressa que la multiplicació de nombres racionals no depèn del representant triat?
\begin{enumerate}
\opcio{a}{\(\left[\frac{ac}{bd}\right]=\left[\frac{a'c}{b'd}\right]\).}
\opcio{b}{\(\left[\frac{a+c}{b+d}\right]=\left[\frac{a'+c}{b'+d}\right]\).}
\opcio{c}{\(\left[\frac{a}{b}\right]=\left[\frac{a'}{d}\right]\).}
\opcio{d}{La multiplicació no està ben definida en \(\Q\).}
\end{enumerate}
\feedback

\pregunta Per situar el nombre racional
\[
	-\frac{11}{4}
\]
a la recta numèrica mitjançant el teorema de Tales, quin procediment és correcte?
\begin{enumerate}
\opcio{a}{Escrivim \(-\frac{11}{4}=-2+\frac{3}{4}\) i dividim \([-2,-1]\) en quatre parts iguals.}
\opcio{b}{Dividim directament el segment \([0,-1]\) en onze parts iguals.}
\opcio{c}{Escrivim \(-\frac{11}{4}=-3+\frac{1}{4}\), dividim \([-3,-2]\) en quatre parts iguals i prenem el primer punt des de \(-3\) cap a \(-2\).}
\opcio{d}{Escrivim \(-\frac{11}{4}=-4+\frac{3}{4}\) i dividim \([-4,-3]\) en tres parts iguals.}
\end{enumerate}
\feedback

\pregunta Quin dels nombres següents està estrictament entre
\[
	\frac{5}{7}
	\qquad\text{i}\qquad
	\frac{3}{4}?
\]
\begin{enumerate}
\opcio{a}{\(\frac{29}{42}\).}
\opcio{b}{\(\frac{41}{56}\).}
\opcio{c}{\(\frac{7}{10}\).}
\opcio{d}{\(\frac{3}{5}\).}
\end{enumerate}
\feedback

\pregunta La fracció irreductible
\[
	\frac{17}{120}
\]
té una expressió decimal:
\begin{enumerate}
\opcio{a}{Finita, perquè \(120\) és múltiple de \(10\).}
\opcio{b}{Periòdica pura, perquè \(120\) no és potència de \(10\).}
\opcio{c}{Entera, perquè \(17\) i \(120\) són primers entre si.}
\opcio{d}{Periòdica mixta, perquè \(120=2^3\cdot3\cdot5\).}
\end{enumerate}
\feedback

\pregunta El resultat de
\[
	\left(\frac{3}{4}-\frac{5}{6}\right):\left(-\frac{1}{3}\right)
	+
	\left(-\frac{2}{5}\right)^2
\]
és:
\begin{enumerate}
\opcio{a}{\(\frac{29}{100}\).}
\opcio{b}{\(\frac{41}{100}\).}
\opcio{c}{\(\frac{1}{20}\).}
\opcio{d}{\(\frac{9}{100}\).}
\end{enumerate}
\feedback

\pregunta El resultat de
\[
	\sqrt{\frac{49}{81}}
	-
	\left(-\frac{2}{3}\right)^2
	+
	\sqrt[3]{-\frac{8}{27}}
\]
és:
\begin{enumerate}
\opcio{a}{\(\frac{1}{3}\).}
\opcio{b}{\(-\frac{1}{3}\).}
\opcio{c}{\(-\frac{7}{9}\).}
\opcio{d}{\(\frac{1}{9}\).}
\end{enumerate}
\feedback

\pregunta Si \(x\) és la quarta proporcional de
\[
	\frac{3}{5},
	\qquad
	\frac{2}{3}
	\qquad\text{i}\qquad
	\frac{4}{9},
\]
aleshores
\[
	\frac{\frac{3}{5}}{\frac{2}{3}}
	=
	\frac{\frac{4}{9}}{x}.
\]
Quin valor té \(x\)?
\begin{enumerate}
\opcio{a}{\(\frac{8}{27}\).}
\opcio{b}{\(\frac{40}{81}\).}
\opcio{c}{\(\frac{6}{5}\).}
\opcio{d}{\(\frac{81}{40}\).}
\end{enumerate}
\feedback

\pregunta La mitjana proporcional positiva de
\[
	\frac{9}{4}
	\qquad\text{i}\qquad
	\frac{16}{9}
\]
és:
\begin{enumerate}
\opcio{a}{\(\frac{13}{6}\).}
\opcio{b}{\(4\).}
\opcio{c}{\(2\).}
\opcio{d}{\(\frac{1}{2}\).}
\end{enumerate}
\feedback

\pregunta Considerem la taula
\[
\begin{array}{c|ccc}
x & \frac{1}{2} & \frac{3}{4} & \frac{3}{2} \\
\hline
y & 12 & 8 & 4
\end{array}
\]
La relació entre \(x\) i \(y\) és:
\begin{enumerate}
\opcio{a}{De proporcionalitat directa, perquè \(y\) disminueix quan \(x\) augmenta.}
\opcio{b}{De proporcionalitat inversa, perquè el producte \(xy\) és constant.}
\opcio{c}{De proporcionalitat directa, perquè el quocient \(\frac{y}{x}\) és constant.}
\opcio{d}{Cap relació de proporcionalitat.}
\end{enumerate}
\feedback

\pregunta Cinc treballadors, treballant \(6\) hores al dia, acaben una feina en \(12\) dies. Si hi treballen \(8\) treballadors, \(5\) hores al dia, quants dies trigaran a acabar la mateixa feina, suposant que tots treballen al mateix ritme?
\begin{enumerate}
\opcio{a}{\(8\) dies.}
\opcio{b}{\(9\) dies.}
\opcio{c}{\(10\) dies.}
\opcio{d}{\(16\) dies.}
\end{enumerate}
\feedback

\pregunta Tres màquines fan una mateixa tasca en \(4\), \(5\) i \(10\) hores, respectivament. Es vol repartir una prima de \(660\) euros premiant més la màquina més ràpida. Quin repartiment correspon a aquest criteri?
\begin{enumerate}
\opcio{a}{\(220\), \(220\), \(220\).}
\opcio{b}{\(120\), \(240\), \(300\).}
\opcio{c}{\(300\), \(240\), \(120\).}
\opcio{d}{\(250\), \(220\), \(190\).}
\end{enumerate}
\feedback

\pregunta Una funció de proporcionalitat inversa compleix
\[
	f\left(-\frac{3}{2}\right)=4.
\]
Quin és el valor de
\[
	f\left(\frac{9}{4}\right)?
\]
\begin{enumerate}
\opcio{a}{\(-\frac{8}{3}\).}
\opcio{b}{\(\frac{8}{3}\).}
\opcio{c}{\(-\frac{27}{8}\).}
\opcio{d}{\(-\frac{2}{3}\).}
\end{enumerate}
\feedback
\end{enumerate}
